\section{模块化无翼/有翼垂起形态的VTOL分析与仿真}
\label{sec:vtol}

本节探讨纵列双发正交单轴矢量推力平台在不同升力形态下实现垂直起降（VTOL）的局部可行性。
这里的“可行性”仅指在给定集总参数、推力裕度和姿态控制律下的局部动力学响应，
不等同于已经完成无翼机体的气动外形、结构强度、散热、能源和安全性设计。
其中，tail-sitter是本文采用的一个具体姿态方案：飞行器从水平姿态通过俯仰$-90^\circ$变为竖直姿态（机头朝天），
两电机推力在惯性系中提供主要竖直升力。
在无翼形态下，气动升力不再是悬停的必要条件，推进器承担主要升力和姿态控制；
在有翼形态下，机翼可在前飞阶段逐步承担升力，但垂起段仍主要依靠推进器。
第\ref{sec:allocation}节和第\ref{sec:propulsion}节推导的推力矢量映射和推进模型对两种形态均适用，
但质量、惯量、推重比和气动参数应按式(\ref{eq:configuration_params})分别标定。

\subsection{Tail-sitter构型的控制特性}

在tail-sitter悬停姿态（项目兼容显示角$\theta=-90^\circ$，Web机体系$+x_b$指向NED的$-z$方向）下，
构型的控制特性与有翼前飞形态有本质变化；但其核心推进器、摆座和差速通道仍保持不变：

\begin{itemize}
    \item \textbf{通道语义按悬停构型重读}：在推进模型系中，前电机摆座绕$z'$轴、尾电机摆座绕$y'$轴；经式(\ref{eq:frame_permutation})和实机执行器标定后，任务语义为\textbf{上摆=滚转主控、下摆=俯仰主控}。这里必须区分物理摆轴与任务轴名称，不能把$z'$直接写成固件$z_b$（通道命名见第\ref{sec:axis_semantics}节）。
    \item \textbf{差速=偏航主控}：前后电机反扭矩差产生绕推力轴（机身纵轴）的力矩，悬停构型下即偏航——该机制不依赖重力方向，VTOL模式下完全保留。
\item \textbf{航向锁（ratchet hold，实机 2026-08-09）}：悬停航向不做姿态回中（摇杆恒为角速度指令），
回中时以\textbf{航向短弧误差} $\Delta\psi\in(-\pi,\pi]$ 经比例增益（$k_p^{\text{yaw}}=5.0$，限幅 60°/s）
生成角速率指令锁住当时航向——"松杆即停"由航向锁替代（存档版为自由漂移），
解锁瞬间参考重置为当前航向防起飞猛转。
该实现绕开了第\ref{sec:quat_control}章的四元数航向回中（悬停万向锁附近欧拉表示退化），
是工程上规避欧拉奇异 + 满足"摇杆恒角速度"操作习惯的折中。
    \item \textbf{悬停配平}：忽略气动力后，悬停仅需推力与重力平衡：$T_f+T_t = mg$。在$\delta_f=\delta_t=0$、$\Delta\omega=0$（零差速）的对称状态下，解析解得基准转速
    \begin{equation}
        \omega_{0,\text{hover}} = \sqrt{\frac{mg}{2k_T}}
        \label{eq:hover_omega0}
    \end{equation}
    代入VTOL参数集（$m=0.7$\,kg，$k_T=1.04\times10^{-5}$\,N$\cdot$s$^2$，v0.2.0）得$\omega_{0,\text{hover}} \approx 574.0$\,rad/s，对应油门约$49.9\%$。
\end{itemize}

\subsection{四元数控制在VTOL中的必要性}

Tail-sitter悬停时俯仰角$\theta=-90^\circ$恰为3-2-1欧拉角的万向锁奇异点。
按第\ref{sec:quat_control}节当前到期望的误差约定，Web期望悬停四元数为
$\mathbf q_{\text{hover}}=[0.7071,0,+0.7071,0]^T$；该符号由
\texttt{Q\_HOVER}和“$+x_b\to-z_{\mathrm{NED}}$”回归断言锁定。
四元数本身在此姿态正则，但闭环仍需处理$180^\circ$短弧选择、执行器饱和和任务轴置换。

从水平姿态到悬停姿态的参考可采用slerp球面线性插值（式\ref{eq:slerp}）。
在端点先执行短弧符号选择且插值后保持归一化时，中间参考是有效单位四元数并沿所选大圆弧行进；
若还要求起止角速度连续，则需对插值参数增加缓入缓出时间标度。

\subsection{仿真验证}

本节四张图生成于当前参数源与固件轴置换之前，使用历史控制脚本和$\pm25^\circ$模型限幅。
为保留研究过程，图与原数值继续列出，但均标为历史脚本输出；在按当前参数、四元数方向和
$\pm15^\circ$固件限幅重跑前，不作为当前v0.2.0或实机性能证据。

\subsubsection{垂直起飞}

图\ref{fig:vtol_takeoff}记录历史脚本的3秒slerp过渡。该脚本报告误差范数峰值约0.02、
过渡后回落以及10\,s末高度约28.5\,m。由于地面接触、推力限幅、坐标和控制器均未按当前
管线重建，这些数值只描述旧仿真轨迹，不构成起飞高度或跟踪性能预测。

\begin{figure}[H]
\centering
\includegraphics[width=0.88\textwidth]{fig-sim/vtol_takeoff.pdf}
\caption{历史VTOL脚本的3\,s slerp轨迹：高度、俯仰角及四元数误差范数。
图中LQR和$25^\circ$限幅均为旧口径。}
\label{fig:vtol_takeoff}
\end{figure}

三种历史控制器的输出见图\ref{fig:vtol_takeoff_compare}和表\ref{tab:ctrl_compare}。
旧脚本分别给出约$23^\circ$、$6^\circ$和$2^\circ$的回摆，但三者并未共享当前固件的
误差方向、轴置换、限幅与控制链，故只保留为研究过程记录，不解释为性能梯度。

\begin{figure}[H]
\centering
\includegraphics[width=0.88\textwidth]{fig-sim/vtol_takeoff_compare.pdf}
\caption{历史3\,s slerp脚本中的俯仰角与四元数误差范数。
曲线仅能在该旧模型内部比较。}
\label{fig:vtol_takeoff_compare}
\end{figure}

\subsubsection{悬停保持}

图\ref{fig:vtol_hover}展示历史脚本在解析悬停转速下的定点保持结果。
该无扰动、参数匹配场景在8\,s内未显示高度漂移，说明数值模型可保持所设平衡点；
它不能证明误差四元数控制在扰动、饱和或参数失配下的稳定性。四元数状态在欧拉角
$\theta=-90^\circ$处仍是正则的，这是表示法性质而非该单次仿真的性能结论。

\begin{figure}[H]
\centering
\includegraphics[width=0.88\textwidth]{fig-sim/vtol_hover_hold.pdf}
\caption{历史无扰动悬停脚本中的俯仰角和高度。}
\label{fig:vtol_hover}
\end{figure}

\subsubsection{三种控制器的悬停抗扰对比}

历史脚本还施加了$3^\circ$初始俯仰偏差，三条误差曲线见
图\ref{fig:vtol_ctrl_compare}。原输出中SAS出现回摆，INDI近似单调，LQR在5\,s末的
误差范数最小；由于控制器与当前实现不一致，这些形态不能用于解释INDI鲁棒性或LQR优势。

\begin{figure}[H]
\centering
\includegraphics[width=0.82\textwidth]{fig-sim/vtol_ctrl_compare.pdf}
\caption{历史$3^\circ$初始俯仰偏差脚本的三控制器误差曲线。}
\label{fig:vtol_ctrl_compare}
\end{figure}

上述精确峰值、收敛时间和高度只对应历史脚本。由于图资产早于当前模型与固件，
本章不再把单组无风、无参数失配试验解释为INDI鲁棒性、LQR普适优势或三种方法的性能梯度。
正式VTOL对比需由当前参数源重跑，并自动记录命令、提交、限幅、控制器定义与输出摘要。

\subsection{VTOL模式的已知局限}

\begin{enumerate}[label=(\arabic*)]
    \item \textbf{悬停工作点与推力裕度}：$49.9\%$表示指令转速占$\omega_{\max}$的比例，
          不是“仅余50\%推力”。按$T\propto\omega^2$，MODEL-DEFAULT最大推重比约为
          $(1150/574)^2\approx4.0$；这表示模型具有较大名义推力裕度，但$k_T$、
          电池电压与带桨极限尚未台架标定，不能据此确认实机裕度；
    \item \textbf{有翼模块在垂起段的利用率低}：在垂起模式下，机翼和舵面通常不承担主要升力与姿态控制，
          但这不是无翼形态的缺点，而是有翼模块为前飞效率换取的质量和阻力代价。
    \item \textbf{形态/姿态过渡建模缺失}：当前仿真未建模有翼/无翼参数切换及slerp过渡中的非定常气动效应，
          机翼在过渡段可能产生显著的阻力和力矩。
    \item \textbf{单轴摆座的平动力方向约束}：每台推进器只能在各自摆动平面内改变推力方向，
          因而两台电机不能像双自由度万向节那样独立指定任意横向力矢量。过渡段中的惯性系
          竖直分量并非简单“缺失”，而是由机体姿态、总推力和两个正交摆角共同决定；真正的限制是
          平动力目标与姿态力矩目标相互耦合，必须通过可达集和饱和约束分析，而不能仅凭单个
          机体系分量判定某一阶段不可实现。
\end{enumerate}

\subsection{不同形态的统一控制架构}

将有翼前飞、无翼垂起和有翼垂起统一于四元数姿态控制和在线效能矩阵框架下：
高层根据形态和任务切换期望四元数$\mathbf{q}_{\text{des}}$、推力指令、气动模型参数和控制权限调度；
核心控制律（例如第\ref{sec:cascade_chapter}节的误差四元数外环与$\mathbf{B}_{\text{true}}$内环）保持不变。
双推进器和摆座提供跨形态的基础控制，机翼在有翼前飞段增加气动升力和舵面控制，
无翼形态则由推进器承担主要平动与姿态控制。形态切换必须同步处理质量/惯量变化、配平变化、
控制权重变化和执行器限幅，不能只切换一个模式标志。
